% Study 1: effect sizes.
% Written by tools/make_tables.py from the pipeline table(s) below (paths relative to
% Analysis/R Analysis Pipeline, each with the SHA-256 it was read with); do not edit by hand.
%   Study 1/tables/study1_bootstrap_d.csv  a223905d69ebbaaa35cb1f512a1a1cdefca4c5678b8057acd613e38876f048ca
\begin{table}[htbp]
\centering
\small
\caption[Study 1: effect sizes]{\textbf{Study~1: effect sizes of the label's penalty.} The registered effect sizes (registered secondary; no verdict). Averaging each person's ten ratings first removes much of the variation, so $d$ on participants' mean ratings is larger than the model-based $d$ by construction and cannot be read against the usual small, medium and large benchmarks. The model-based $d$ is the one to compare with other studies.}\label{tab:supp_bootstrap_d}
\setlength{\tabcolsep}{4pt}
\begin{tabular}{@{}lccc@{}}
\toprule
Rating & $d$ on participants' means & 95\% CI & Model-based $d$\\
\midrule
Creativity & 0.66 & 0.30 to 1.06 & 0.41\\
Effort & 1.15 & 0.77 to 1.64 & 0.80\\
\bottomrule
\end{tabular}
\par\smallskip
\begin{minipage}{\textwidth}\footnotesize\raggedright Cohen's $d$ here is the Human-label mean minus the AI-label mean, divided by the pooled standard deviation of participants' mean ratings (56 Human-label and 55 AI-label participants). Its 95\% confidence interval (CI) comes from 10{,}000 resamples of participants drawn within each label group. The model-based $d$ divides the H1 contrast by the square root of the primary model's summed variances (between participants, between images, and the rest), which puts it on the scale of a single rating.\par\end{minipage}
\end{table}
